%%%%%%%%%%  scan idx 292  |  folio 279  %%%%%%%%%%
\SGAPrintPage{279}%
The last two categories are full subcategories of the category \(\Mod({}_A X)\) of left \(A\)-Modules, and the first two are full triangulated subcategories of \(D({}_A X)\) (\(= D(\Mod({}_A X))\)).

\sgaStableTarget{sga6.e04.def.2-2}\textbf{Definition 2.2.} \uline{Let} \((X,A)\) \uline{be a ringed topos. Unless otherwise stated, we shall put}
\[
 \begin{gathered}
  K_\bullet({}_A X)=k(D^b({}_A X)_{\coh}),\\
  K^{\bullet}({}_A X)=k(D({}_A X)_{\uparf}),
 \end{gathered}
\]
\uline{where} \(k(\ )\) \uline{is defined in} (1.1)\uline{, and}
\[
 \begin{gathered}
  K_\bullet({}_A X)_{\naif}=k(\Coh({}_A X)),\\
  K^{\bullet}({}_A X)_{\naif}=k(\Loclib({}_A X)),
 \end{gathered}
\]
\uline{where} \(k(\ )\) \uline{is defined by} (1.5.2)\uline{, the ambient abelian category being} \(\Mod({}_A X)\)\SGAAuthorFootnote{1}{We shall allow ourselves to omit the index \(A\) when no confusion can result. Sometimes we shall also drop the index ``naïf'', but shall indicate this explicitly.}\uline{. We shall say that} \(K_\bullet({}_A X)\) \uline{(resp.} \(K^{\bullet}({}_A X)\)\uline{) is the Grothendieck group of pseudo-coherent (resp. perfect) complexes of left} \(A\)\uline{-Modules, and that} \(K_\bullet({}_A X)_{\naif}\) \uline{(resp.} \(K^{\bullet}({}_A X)_{\naif}\)\uline{) is the Grothendieck group of coherent (resp. locally free) left} \(A\)\uline{-Modules.}

\sgaStableTarget{sga6.e04.rem.2-3}\textbf{Remarks 2.3.}
\begin{enumerate}[label=\alph*)]
\item One can define variants of \(K_\bullet(\ )\) and \(K^{\bullet}(\ )\) by putting
\[
 K_\bullet^{\loc}({}_A X)=k(D^{\loc b}({}_A X)_{\coh})
\]
(where ``\(\loc b\)'' means ``with locally bounded cohomology''), and
\[
 K^{\bullet}_{\loc}({}_A X)=k(D({}_A X)_{\mathrm{parf}}).
\]
These groups coincide with those defined in (2.2) when the final object of \(X\) is quasi-compact, but in general are less suited to intersection calculations.

%%%%%%%%%%  scan idx 293  |  folio 280  %%%%%%%%%%
\item\SGAPrintPage{280}%
The group \(K_\bullet({}_A X)_{\naif}\) is practically of interest only when \(A\) is coherent (otherwise there may well be very few coherent sheaves \ldots).
\item Considering the \(k\) of triangulated categories that are ``too large'' is of no interest. Thus, from (1.2 b)) one easily deduces the relations
\[
 k(D({}_A X))=k(D^-({}_A X))=k(D^+({}_A X))=k(D^b({}_A X))=k(D^-({}_A X)_{\coh})=0
\]
(for the last equality, observe that \(L\in\Ob D^-({}_A X)_{\coh}\) implies \(\coprod_{i\geq0}L[2i]\in\Ob D^-({}_A X)_{\coh}\)).
\end{enumerate}

\sgaStableTarget{sga6.e04.para.2-4}2.4. \uline{Case of a coherent sheaf of rings}. Let \((X,A)\) be a ringed topos, with \(A\) supposed \uline{coherent} as a left \(A\)-Module (I 3.1 and 3.6). Every coherent left \(A\)-Module is a pseudo-coherent complex (I 3.5), so we have a fully faithful functor
\sgaStableTarget{sga6.e04.eq.2-4-1}\begin{equation}\tag{2.4.1}
 \Coh({}_A X)\hookrightarrow D^b({}_A X)_{\coh}.
\end{equation}
Now (I 3.4) \(\Coh({}_A X)\) is a thick subcategory of \(\Mod({}_A X)\), and, by (I 3.5), \(D^b({}_A X)_{\coh}\) is the full subcategory of \(D^b({}_A X)\) formed by the objects \(L\) such that \(H^i(L)\in\Ob\Coh({}_A X)\) for every \(i\). Consequently, it follows from (1.6) that \uline{the homomorphism}
\sgaStableTarget{sga6.e04.eq.2-4-2}\begin{equation}\tag{2.4.2}
 K_\bullet({}_A X)_{\naif}\longrightarrow K_\bullet({}_A X)
\end{equation}
\uline{defined by} (2.4.1) \uline{is an isomorphism}.

\sgaStableTarget{sga6.e04.para.2-5}2.5. \uline{The homomorphism} \(K^{\bullet}\to K_\bullet\). Let \((X,A)\) be a ringed topos. The inclusion functor \(D({}_A X)_{\uparf}\hookrightarrow D^b({}_A X)_{\coh}\) defines a homomorphism
\sgaStableTarget{sga6.e04.eq.2-5-1}\begin{equation}\tag{2.5.1}
 K^{\bullet}({}_A X)\longrightarrow K_\bullet({}_A X).
\end{equation}
This is in general neither injective nor surjective. However, make the following hypotheses:
\begin{enumerate}[label=(\roman*)]

%%%%%%%%%%  scan idx 294  |  folio 281  %%%%%%%%%%
\item\SGAPrintPage{281}%
the final object of \(X\) is quasi-compact;
\item \(X\) has enough points, and, for every point \(x\) of \(X\), every left \(A_x\)-module has finite tor-dimension.
\end{enumerate}
Then, by (I 5.10), \(D({}_A X)_{\parf}=D^b({}_A X)_{\coh}\), and the homomorphism (2.5.1) is an isomorphism.

Hypotheses (i) and (ii) are satisfied, for example, by a compact space ringed by regular local rings.

\sgaStableTarget{sga6.e04.para.2-6}2.6. \uline{Pairings between} \(K^{\bullet}\) \uline{and} \(K^{\bullet}\) (\uline{noncommutative case}). Let
\[
 f:(X',A')\longrightarrow (X,A)
\]
be a morphism of ringed topoi. Denote by \(\Mod({}_{A'} X'_A)\) the category of bi-Modules, left over \(A'\) and right over \(f^{-1}(A)\), and by \(\Loclib_{A'}({}_{A'} X'_A)\) the full subcategory formed by the objects belonging to \(\Loclib({}_{A'} X')\). By analogy with (2.2), we shall put
\[
 K^{\bullet}_{\naif}({}_{A'} X'_A)_{A'}=k(\Loclib_{A'}({}_{A'} X'_A)).
\]
The tensor product
\[
 \otimes_A:\Mod({}_{A'} X'_A)\times\Mod({}_A X)\longrightarrow \Mod({}_{A'} X')
\]
\[
 (E',E)\longmapsto E'\otimes_{f^{-1}(A)}f^{-1}(E)
\]
induces a functor
\sgaStableTarget{sga6.e04.eq.2-6-1}\begin{equation}\tag{2.6.1}
 \Loclib_{A'}({}_{A'} X'_A)\times\Loclib({}_A X)\longrightarrow\Loclib({}_{A'} X').
\end{equation}
The map \((E',E)\mapsto\cl(E'\otimes_A E)\) deduced from (2.6.1) is biadditive in the sense of (1.5), and therefore defines a homomorphism
\sgaStableTarget{sga6.e04.eq.2-6-2}\begin{equation}\tag{2.6.2}
 K^{\bullet}_{\naif}({}_{A'} X'_A)_{A'}\otimes K^{\bullet}_{\naif}({}_A X)\longrightarrow K^{\bullet}_{\naif}({}_{A'} X').
\end{equation}
Denote by \(D({}_{A'} X'_A)\) the derived category of \(\Mod({}_{A'} X'_A)\), and by \(D({}_{A'} X'_A)_{A'\text{-}\parf}\) the full subcategory formed by the complexes of finite perfect amplitude rel. to \(A'\) (i.e. as objects of \(D({}_{A'} X')\)), and put, as in (2.2),
\[
 K^{\bullet}({}_{A'} X'_A)_{A'}=k(D({}_{A'} X'_A)_{A'\text{-}\parf}).
\]

%%%%%%%%%%  scan idx 295  |  folio 282  %%%%%%%%%%
\SGAPrintPage{282}%
The inclusion functors plainly define homomorphisms
\sgaStableTarget{sga6.e04.eq.2-6-3-1}\sgaStableTarget{sga6.e04.eq.2-6-3-2}\sgaStableTarget{sga6.e04.eq.2-6-3-3}\begin{align}
 K^{\bullet}_{\naif}({}_{A'} X'_A)_{A'}&\longrightarrow K^{\bullet}({}_{A'} X'_A)_{A'}, \tag{2.6.3.1}\\
 K^{\bullet}_{\naif}({}_A X)&\longrightarrow K^{\bullet}({}_A X), \tag{2.6.3.2}\\
 K^{\bullet}_{\naif}({}_{A'} X')&\longrightarrow K^{\bullet}({}_{A'} X'). \tag{2.6.3.3}
\end{align}
On the other hand, we know (I 4.18, 4.19) that the derived functor of \(\otimes_A\),
\[
 \otimes_A^{\mathrm L}:D^-({}_{A'} X'_A)\times D^-({}_A X)\longrightarrow D^-({}_{A'} X'),
\]
induces an (exact) functor
\sgaStableTarget{sga6.e04.eq.2-6-4}\begin{equation}\tag{2.6.4}
 D({}_{A'} X'_A)_{A'\text{-}\parf}\times D({}_A X)_{\parf}
 \longrightarrow D({}_{A'} X')_{\parf}.
\end{equation}
The latter defines a homomorphism
\sgaStableTarget{sga6.e04.eq.2-6-5}\begin{equation}\tag{2.6.5}
 K^{\bullet}({}_{A'} X'_A)_{A'}\otimes K^{\bullet}({}_A X)\longrightarrow K^{\bullet}({}_{A'} X').
\end{equation}
Since the functor (2.6.1) is induced by (2.6.4), the homomorphisms (2.6.2) and (2.6.5) are compatible with the homomorphisms (2.6.3.1), (2.6.3.2), and (2.6.3.3); in other words, we have a commutative diagram
\sgaStableTarget{sga6.e04.eq.2-6-6}\[
\begin{tikzcd}[column sep=tiny]
K^{\bullet}_{\naif}({}_{A'} X'_A)_{A'}\otimes K^{\bullet}_{\naif}({}_A X)\arrow[r]\arrow[d]
&K^{\bullet}_{\naif}({}_{A'} X')\arrow[d]\\
K^{\bullet}({}_{A'} X'_A)_{A'}\otimes K^{\bullet}({}_A X)\arrow[r]
&K^{\bullet}({}_{A'} X'),
\end{tikzcd}
\tag{2.6.6}
\]
where the upper and lower horizontal arrows are respectively (2.6.2) and (2.6.5), and the vertical arrows are defined by (2.6.3.1), (2.6.3.2), and (2.6.3.3).

%%%%%%%%%%  scan idx 296  |  folio 283  %%%%%%%%%%
\SGAPrintPage{283}%
\sgaStableTarget{sga6.e04.para.2-7}2.7. \uline{Contravariant law and ring structure on} \(K^{\bullet}\) We shall spell out some important special cases of the pairings (2.6) (their direct study is moreover immediate).
\begin{enumerate}[label=\alph*)]
\item Let \(f:(X',A')\longrightarrow (X,A)\) be a morphism of ringed topoi. The functor \(Lf^*:D^-({}_A X)\longrightarrow D^-({}_{A'} X')\) sends (I 4.19.1) \(D({}_A X)_{\parf}\) to \(D({}_{A'} X')_{\parf}\) (resp. \(\Loclib({}_A X)\) to \(\Loclib({}_{A'} X')\)), and defines a homomorphism
\sgaStableTarget{sga6.e04.eq.2-7-1}\begin{equation}\tag{2.7.1}
 f^*:K^{\bullet}({}_A X)\longrightarrow K^{\bullet}({}_{A'} X')
\end{equation}
(resp.
\sgaStableTarget{sga6.e04.eq.2-7-2}\begin{equation}\tag{2.7.2}
 f^*:K^{\bullet}({}_A X)_{\naif}\longrightarrow K^{\bullet}({}_{A'} X')_{\naif}
\end{equation}
), which can also be interpreted as multiplication by \(\cl(A')\) in (2.6.5) (resp. (2.6.2)). We have a commutative diagram (a special case of (2.6.6))
\[
\begin{tikzcd}
K^{\bullet}({}_A X)_{\naif}\arrow[r,"f^*"]\arrow[d]
&K^{\bullet}({}_{A'} X')_{\naif}\arrow[d]\\
K^{\bullet}({}_A X)\arrow[r,"f^*"]
&K^{\bullet}({}_{A'} X'),
\end{tikzcd}
\]
where the left and right vertical arrows are respectively (2.6.3.2) and (2.6.3.3).

Let \(g:(X,A)\to(X'',A'')\) be a morphism of ringed topoi. The transitivity of inverse-image functors implies that the homomorphisms of type (2.7.1) (resp. 2.7.2) satisfy the relation
\sgaStableTarget{sga6.e04.eq.2-7-3}\begin{equation}\tag{2.7.3}
 (gf)^*=f^*g^*.
\end{equation}
Thus \(K^{\bullet}({}_A X)\) (resp. \(K^{\bullet}({}_A X)_{\naif}\)) \uline{is a contravariant functor of} \((X,A)\).

\item Let \((X,\underline{O}_X)\) be a commutatively ringed topos. By the associativity and commutativity of the tensor product, \uline{the homomorphism} (2.6.5) (resp. (2.6.2)) (where we put \(A=\underline{O}_X\), \((X',A')=(X,A)\), \(f=\Id\)) \uline{endows} \(K^{\bullet}(X)\) (resp. \(K^{\bullet}(X)_{\naif}\)) \uline{with the structure of a commutative ring with unit}
